% Figure: Fresnel zone plate. Left: concentric opaque and transparent rings whose
% edges fall at radii r_n = sqrt(n lambda f), so that alternate half-period zones
% are blocked and the transmitted contributions add in phase at the focus. Right:
% the diffractive focusing of a plane wave to the primary focal point, with the
% higher diffraction orders forming weaker foci at f/3, f/5, and so on.
\begin{figure}[htbp]
\centering
\begin{tikzpicture}[>=Stealth, scale=1.0]
% left panel: the zone pattern
\begin{scope}[shift={(0,0)}]
% outer disc transparent (white); fill opaque even zones dark
\draw[engblack] (0,0) circle [radius=1.9];
\foreach \n in {1,2,3,4,5,6} {
  \pgfmathsetmacro{\rout}{1.9*sqrt(\n/6)}
  \pgfmathsetmacro{\rin}{1.9*sqrt((\n-1)/6)}
  \ifodd\n
    \fill[engblue!30] (0,0) circle [radius=\rout];
    \fill[white] (0,0) circle [radius=\rin];
  \fi
}
% redraw the odd shaded rings on top of white cores (loop already handles by order)
\foreach \n in {1,3,5} {
  \pgfmathsetmacro{\rout}{1.9*sqrt(\n/6)}
  \pgfmathsetmacro{\rin}{1.9*sqrt((\n-1)/6)}
  \fill[engblue!30] (0,0) circle [radius=\rout];
  \fill[white] (0,0) circle [radius=\rin];
}
\draw[engblack] (0,0) circle [radius=1.9];
\node[font=\scriptsize, anchor=north] at (0,-2.15) {alternate zones blocked};
\draw[<-, enggray] (0,0) -- ({1.9*0.816},0);
\node[enggray, font=\scriptsize, anchor=south] at (1.0,0.05) {$r_n=\sqrt{n\lambda f}$};
\end{scope}
% right panel: focusing
\begin{scope}[shift={(6.2,0)}]
\draw[enggray, dashed] (-2.0,0) -- (3.0,0);
% zone plate as a vertical hatched bar at x=-1.4
\draw[very thick, engblue] (-1.4,-1.5) -- (-1.4,1.5);
\node[engblue, font=\scriptsize, anchor=south] at (-1.4,1.5) {zone plate};
% incoming plane wave arrows
\foreach \y in {-1.1,-0.55,0.55,1.1} {
  \draw[engorange, ->, thick] (-2.0,\y) -- (-1.42,\y);
}
% converging to primary focus at x=1.6
\draw[engamber, thick] (-1.4,1.5) -- (1.6,0);
\draw[engamber, thick] (-1.4,-1.5) -- (1.6,0);
\draw[engamber, thick] (-1.4,0.75) -- (1.6,0);
\draw[engamber, thick] (-1.4,-0.75) -- (1.6,0);
\filldraw[engblack] (1.6,0) circle [radius=1.3pt];
\node[font=\scriptsize, anchor=north west] at (1.6,-0.05) {focus $f$};
% weaker third-order focus at x=-0.4
\filldraw[engred] (-0.4,0) circle [radius=1.0pt];
\node[engred, font=\scriptsize, anchor=south] at (-0.4,0.05) {$f/3$};
\draw[<->, enggray] (-1.4,-1.75) -- (1.6,-1.75);
\node[enggray, font=\scriptsize, anchor=north] at (0.1,-1.78) {$f$};
\end{scope}
\end{tikzpicture}
\caption[Fresnel zone plate]{A Fresnel zone plate focuses by diffraction rather
than refraction. Its ring edges fall at radii $r_n=\sqrt{n\lambda f}$, so that
each transparent zone contributes light that arrives at the axial focus within
half a wavelength of every other transparent zone's contribution, and the
transmitted wavelets add nearly in phase. Blocking the alternate (out-of-phase)
zones removes the cancelling contributions and leaves a real focus at distance
$f=r_1^2/\lambda$. Because the focal length carries $\lambda$ in the denominator,
the zone plate is strongly chromatic, and it forms weaker higher-order foci at
$f/3$, $f/5$, and so on. Zone plates focus X-rays and extreme ultraviolet, where
no refracting lens exists, and pattern the focusing elements of some
compact optical sensors.}
\label{fig:vol04ch11:zone-plate}
\end{figure}
